% =========================================================
% Satisfy-and-Cite Uniqueness
% =========================================================
\paragraph{The Satisfy-and-Cite Pattern.}
\label{sec:satisfy-cite}

There is a proof pattern that does not appear in any of the standard
proof archetypes by name, yet appears constantly --- in every chapter of
abstract algebra and analysis. Most undergraduates have never been taught
it explicitly, which means they either reinvent it poorly each time or
miss it entirely and write longer, less clear proofs.

The pattern is called \textbf{satisfy-and-cite}. Here is its logic.

\begin{tcolorbox}[colback=propbox, colframe=propborder, arc=2pt,
  left=6pt, right=6pt, top=4pt, bottom=4pt,
  title={\small\textbf{The Satisfy-and-Cite Pattern}}, fonttitle=\small\bfseries]
\textbf{Setup.} You want to prove $a = b$, where $b$ is a
\emph{uniquely determined} object: the unique identity, the unique
inverse, the unique limit, the unique fixed point, the unique
element satisfying some algebraic condition.

\medskip
\textbf{Step 1.} Show that $a$ satisfies the defining property of $b$.

\textbf{Step 2.} Cite the previously proved uniqueness theorem: there is
exactly one object satisfying that property.

\textbf{Step 3.} Conclude that $a = b$, since $b$ is the unique such object
and $a$ is one too.
\end{tcolorbox}

\begin{remark}[Why this is not the same as assume-two-and-compare]
The confusion between these two patterns is extremely common, so it is
worth being precise about what each one does.

The \textbf{assume-two-and-compare} pattern \emph{proves} a uniqueness
theorem. You use it when the goal is to establish that some object is
unique. You say: let $x$ and $y$ both satisfy the definition; then show
$x = y$.

The \textbf{satisfy-and-cite} pattern \emph{uses} an already proved
uniqueness theorem. You use it when you already know some object $b$
is unique and you want to identify $a$ with $b$. You show $a$ satisfies
$b$'s defining property, then invoke uniqueness to collapse $a$ to $b$.

Assume-two-and-compare appears in the proof that the identity is unique.
Satisfy-and-cite appears in every proof that uses that fact afterward.
The first creates the tool; the second uses it. Conflating them means
writing a new uniqueness proof every time you want to identify two equal
objects, which is circular at best and wrong at worst.
\end{remark}

\begin{example}[Socks-shoes property: $(ab)^{-1} = b^{-1}a^{-1}$]
We want to show $(ab)^{-1} = b^{-1}a^{-1}$ in a group.

The inverse of $ab$ is, by definition, the unique element $x$
satisfying $(ab)x = e$ and $x(ab) = e$.

We compute directly, using only associativity and the inverse axiom:
\[
(ab)(b^{-1}a^{-1})
= a(bb^{-1})a^{-1}
= a e a^{-1}
= aa^{-1} = e.
\]
Similarly $(b^{-1}a^{-1})(ab) = e$.

So $b^{-1}a^{-1}$ \emph{satisfies the defining property of $(ab)^{-1}$}.
By the uniqueness of inverses (which we have already proved), we
conclude $(ab)^{-1} = b^{-1}a^{-1}$. \qed

\medskip
\noindent
Notice what did not happen: we never assumed two inverses of $ab$ existed
and compared them. We showed one specific element satisfies the inverse
condition, then cited uniqueness to identify it with the named inverse.
\end{example}

\begin{example}[Ring theorem: $a \cdot (-b) = -(ab)$]
We want to show $a(-b) = -(ab)$ in a ring.

The additive inverse $-(ab)$ is, by definition, the unique element $x$
satisfying $ab + x = 0$.

We compute using distributivity:
\[
ab + a(-b) = a(b + (-b)) = a \cdot 0 = 0.
\]
So $a(-b)$ \emph{satisfies the defining property of $-(ab)$}.
By the uniqueness of additive inverses (previously proved), $a(-b) = -(ab)$. \qed

\medskip
\noindent Same pattern: compute, satisfy, cite.
\end{example}

\begin{remark}[The pattern across mathematics]
Once you see satisfy-and-cite, you will see it everywhere. It appears
wherever uniqueness theorems have been proved:

\medskip
\noindent In \textbf{algebra}: every proof of the form ``this element equals
the inverse/identity/zero'' uses it. The four examples above are
$(ab)^{-1} = b^{-1}a^{-1}$, $a \cdot (-b) = -(ab)$, $-(-a) = a$,
and $a \cdot 0 = 0$.

\noindent In \textbf{analysis}: the uniqueness of limits means that to
show $\lim_{n \to \infty} x_n = L$, you verify $x_n$ converges to $L$
by the definition; uniqueness then implies $L$ equals any other
expression you have shown to be the limit of the same sequence.

\noindent In \textbf{linear algebra}: the zero vector is unique; the
additive inverse is unique; a linear transformation is uniquely determined
by its values on a basis. Each uniqueness theorem creates a new application
of satisfy-and-cite.

\medskip
\noindent
Each time a uniqueness theorem is proved, it creates a new proof tool
that can be activated by satisfy-and-cite for every subsequent argument
in that structure. The theorem is not just a fact --- it is a lever.
\end{remark}
